%!TEX root = ../../exa-ma-d7.3.tex
\subsection{Software: pBB}
\label{sec:pBB:software:training}



\begin{table}[H]
    \centering
    { \setlength{\parindent}{0pt}
    \def\arraystretch{1.25}
    \arrayrulecolor{numpexgray}
    {\fontsize{9}{11}\selectfont
    \begin{tabular}{!{\color{numpexgray}\vrule}p{.4\textwidth}!{\color{numpexgray}\vrule}p{.6\textwidth}!{\color{numpexgray}\vrule}}
        \rowcolor{numpexgray}{\rule{0pt}{2.5ex}\color{white}\bf Field} & {\rule{0pt}{2.5ex}\color{white}\bf Details} \\
        \rowcolor{white}\textbf{Consortium} & \begin{tabular}{l}
Université de Lille\\
\end{tabular} \\
        \rowcolor{numpexlightergray}\textbf{Exa-MA Partners} & \begin{tabular}{l}
Inria Lille\\
\end{tabular} \\
        \rowcolor{white}\textbf{Contact Emails} & \begin{tabular}{l}
nouredine.melab@univ-lille.fr\\
\end{tabular} \\
        \rowcolor{numpexlightergray}\textbf{Supported Architectures} & \begin{tabular}{l}
CPU or GPU\\
\end{tabular} \\
        \rowcolor{white}\textbf{Repository} & \begin{tabular}{l}
\href{https://gitlab.inria.fr/jgmys/permutationbb}{https://gitlab.inria.fr/jgmys/permutationbb} \\
\href{https://github.com/Guillaume-Helbecque/P3D-DFS}{https://github.com/Guillaume-Helbecque/P3D-DFS} \\
\end{tabular} \\
    \rowcolor{numpexlightergray}\textbf{License} & \begin{tabular}{l}
OSS: Cecill-*\\
\end{tabular} \\
        \rowcolor{white}\textbf{Training} & \begin{tabular}{l}
Documentation\\
Tutorial \\
\end{tabular} \\
        \bottomrule
    \end{tabular}

    %%\begin{table}[!ht]
%%    \centering
    {
    \setlength{\parindent}{0pt}
    \def\arraystretch{1.25}
    \arrayrulecolor{numpexgray}
    {\fontsize{9}{11}\selectfont
    \begin{tabular}{!{\color{numpexgray}\vrule}p{0.20\textwidth}
                    !{\color{numpexgray}\vrule}p{0.40\textwidth}
                    !{\color{numpexgray}\vrule}p{0.35\textwidth}
                    !{\color{numpexgray}\vrule}}
    \rowcolor{numpexgray}{\rule{0pt}{2.5ex}\color{white}\bf Training Category}
         & {\rule{0pt}{2.5ex}\color{white}\bf Description}
         & {\rule{0pt}{2.5ex}\color{white}\bf URL/Links} \\
\rowcolor{white}Documentation & User manuals, reference guides, and detailed usage instructions. & \begin{tabular}{p{6cm}}
\url{https://gitlab.inria.fr/jgmys/permutationbb/-/blob/master/README.md}\\
\url{https://github.com/Guillaume-Helbecque/P3D-DFS/blob/main/README.md}\\
\end{tabular} \\
\rowcolor{numpexlightergray}Tutorial & Introduction and code demo for newcomers and practitioners. & \url{https://www.youtube.com/watch?v=jMQ-WYRz-AQ} \\
\bottomrule
    \end{tabular}
    }}
%    \caption{pBB Training}
%\end{table}

    }}
    \caption{pBB Information}
\end{table}


\paragraph{Additional Details:}

\begin{description}

\item[\textbf{Documentation}]
\emph{Short Description: User manuals, reference guides, and detailed usage instructions.}

\begin{itemize}
\item \textbf{Content:} Installation, usage, advanced features, benchmarking, test cases, and API
\item \textbf{Format:} Web pages and runnable Markdown-commented examples
\item \textbf{Development:} An ongoing open-source effort in Markdown and GitHub README files, welcoming community contributions
\end{itemize}

\item[\textbf{Tutorial}]
\emph{Short Description: Introduction and code demo for newcomers and practitioners.}

\begin{itemize}
\item \textbf{Content:} Usage and applications of specific features of pBB
\item \textbf{Format:} Youtube video, part of the ChapelCon '24 video series.
\item \textbf{Target Audience:} Operation research community, newcomers to pBB, Chapel users
\end{itemize}

\end{description}

\subsubsection{Next Steps for pBB}
\label{sec:pBB:next-steps}

Below are some guiding questions to help outline the next steps for training with pBB:
\begin{itemize}
    \item What new training content or updates are planned for pBB in the upcoming milestones?
    \item How will the training materials be updated to incorporate the latest exascale techniques and user feedback?
    \item Which new delivery formats (e.g., webinars, interactive sessions, hands-on workshops) should be introduced?
    \item What additional documentation or support tools are needed to help users get started or deepen their expertise?
    \item How can we better integrate training activities with academic and industrial outreach efforts?
\end{itemize}
